% =====================================================================
% Section: Experiments -- Routing interpretability analysis (outline.md §3.6;
%   defines sec:routing_analysis, referenced by 3_experiments_ablation.tex).
% Draft v1 (2026-06-17). All numbers are taken verbatim from the x4 routing
%   diagnostics logged at iter 80k on Urban100 (paper/data/ablation_metrics.csv,
%   ablation_perblock_urban100_80k.csv). Per-block prototype count
%   M=[16,32,64,128,16,32,64,128] (catanet_arch.py:516); entropy/usage/temperature
%   definitions in catanet_model.py:323-338. Figure: paper/figures/fig6_usage_entropy.tex.
% Claims follow plan/review/method-experiment-traceability.md: usage balancing (C4)
%   is supported by the logged distributions; the per-token x^s histograms and
%   cluster maps require model inference and are reported as qualitative analysis
%   (the toolchain for those visualizations is set up separately). No PSNR claims here.
% =====================================================================
\subsection{Routing interpretability analysis}
\label{sec:routing_analysis}

To inspect DPR's routing behavior independently of its effect on PSNR, we analyze
the diagnostics logged during validation at $\times4$
(iter $80$k, Urban100). For each of the $L{=}8$ TAB blocks we record the per-token
routing confidence $x^{\mathrm{s}}=\max_m S_{n,m}$ (Eq.~\eqref{eq:belong}), the
number of \emph{active} prototypes (slots receiving at least one token), the
normalized usage entropy $\mathcal{H}(u)$ (Eq.~\eqref{eq:balance}), and the learned
temperature $\tau$ (Eq.~\eqref{eq:scores}). The blocks use
$M=[16,32,64,128,16,32,64,128]$ prototypes respectively.

\paragraph{Balanced prototype usage (C4).}
Fig.~\ref{fig:usage_entropy} reports the normalized usage entropy per block with
and without the balancing regularizer. With balancing enabled (the full model) the
entropy is higher in \emph{all eight} blocks, raising the block-mean from $0.6227$
to $0.6865$; the effect is largest in the high-frequency-oriented block $b_4$
($0.7088\!\rightarrow\!0.8466$). Consistently, the average number of active
prototypes increases from $47.5$ to $52.3$ across blocks, meaning that fewer slots are
left unused. Because a learnable temperature can otherwise sharpen the assignment
distribution and concentrate tokens on a handful of prototypes, these two
measurements together show that the regularizer keeps prototype usage spread out
and prevents slot collapse, matching the stabilizing role claimed for C4.

\paragraph{Bounded routing temperature (C4).}
The learnable temperature is clamped by $\tau_{\max}{=}10$ for stability
(Eq.~\eqref{eq:scores}). In the trained full model it settles in a moderate range
across blocks (per-block $\tau$ from $5.65$ to $6.51$, block-mean $6.06$), staying
well below the clamp rather than saturating it. The temperature therefore sharpens
the routing scores enough to separate prototypes while the usage regularizer keeps
the assignment from collapsing---the two C4 components act as intended without
hitting their safety bound.

\paragraph{Confidence-aware routing signal (C3).}
The per-token confidence $x^{\mathrm{s}}$ that drives the ordering key
(Eq.~\eqref{eq:sortkey}), the IASA gate (Eq.~\eqref{eq:iasa_gate}) and the soft
fallback (Eq.~\eqref{eq:soft_fallback}) is a non-degenerate signal in every block.
Fig.~\ref{fig:xscore_hist} shows its distribution accumulated over \emph{all} test
images of Urban100 and Manga109 with the final model: in each block the mass lies
clearly to the right of the uniform-assignment floor $1/M$ (block-mean
$x^{\mathrm{s}}\approx 1.5$--$1.9\times 1/M$, e.g.\ $0.109$ vs.\ $0.0625$ for $b_0$
and $0.012$ vs.\ $0.0078$ for $b_3$ on Urban100) yet stays far from $1$. The
softmax assignment is therefore informative but soft rather than collapsed onto a
single prototype, and its mode shifts toward the lower floor as $M$ grows. Reshaping the
original-order labels $b_n=\arg\max_m S_{n,m}$ to $H\times W$ yields the routing
cluster maps in Fig.~\ref{fig:cluster_maps}: the partitions track image content
(repetitive building facades, high-contrast strokes) instead of a fixed spatial
grid and become finer in the deeper, larger-$M$ blocks. Both views are produced by
a single forward pass with the cached routing state
(\texttt{last\_routing\_map}, \texttt{last\_x\_scores}) and together confirm that
the confidence used by the C3 components is meaningful.

\paragraph{Summary.}
The logged diagnostics provide mechanism-level evidence for the routing design:
the balancing regularizer produces a block-consistent increase in usage
entropy and active-prototype count (C4), the learnable temperature stays in a
stable, unsaturated range (C4), and the retained confidence signal is informative
across all blocks (C3). These observations characterize the behavior of DPR
independently of the reconstruction metrics discussed in
Sec.~\ref{sec:ablation}.
